\documentclass[10pt,a4paper]{amsart}
\usepackage{amsmath,amssymb,amsthm,graphicx}
\usepackage{verbatim}
\newenvironment{reference}{\comment}{\endcomment}
%\newenvironment{reference}{}{}
\newenvironment{slogan}{\comment}{\endcomment}
\newenvironment{history}{\comment}{\endcomment}

% For commutative diagrams we use Xy-pic
\usepackage[all]{xy}

% We use 2cell for 2-commutative diagrams.
\xyoption{2cell}
% Theorem environments.
%
\theoremstyle{plain}
\newtheorem{theorem}[subsection]{Theorem}
\newtheorem{proposition}[subsection]{Proposition}
\newtheorem{lemma}[subsection]{Lemma}

\theoremstyle{definition}
\newtheorem{definition}[subsection]{Definition}
\newtheorem{example}[subsection]{Example}
\newtheorem{exercise}[subsection]{Exercise}
\newtheorem{situation}[subsection]{Situation}

\theoremstyle{remark}
\newtheorem{remark}[subsection]{Remark}
\newtheorem{remarks}[subsection]{Remarks}

\numberwithin{equation}{subsection}

% Macros
%
\def\lim{\mathop{\mathrm{lim}}\nolimits}
\def\colim{\mathop{\mathrm{colim}}\nolimits}
\def\Spec{\mathop{\mathrm{Spec}}}
\def\Hom{\mathop{\mathrm{Hom}}\nolimits}
\def\Ext{\mathop{\mathrm{Ext}}\nolimits}
\def\SheafHom{\mathop{\mathcal{H}\!\mathit{om}}\nolimits}
\def\SheafExt{\mathop{\mathcal{E}\!\mathit{xt}}\nolimits}
\def\Sch{\mathit{Sch}}
\def\Mor{\mathop{\mathrm{Mor}}\nolimits}
\def\Ob{\mathop{\mathrm{Ob}}\nolimits}
\def\Sh{\mathop{\mathit{Sh}}\nolimits}
\def\NL{\mathop{N\!L}\nolimits}
\def\CH{\mathop{\mathrm{CH}}\nolimits}
\def\proetale{{pro\text{-}\acute{e}tale}}
\def\etale{{\acute{e}tale}}
\def\QCoh{\mathit{QCoh}}
\def\Ker{\mathop{\mathrm{Ker}}}
\def\Im{\mathop{\mathrm{Im}}}
\def\Coker{\mathop{\mathrm{Coker}}}
\def\Coim{\mathop{\mathrm{Coim}}}

% Boxtimes
%
\DeclareMathSymbol{\boxtimes}{\mathbin}{AMSa}{"02}

%
% Macros for moduli stacks/spaces
%
\def\QCohstack{\mathcal{QC}\!\mathit{oh}}
\def\Cohstack{\mathcal{C}\!\mathit{oh}}
\def\Spacesstack{\mathcal{S}\!\mathit{paces}}
\def\Quotfunctor{\mathrm{Quot}}
\def\Hilbfunctor{\mathrm{Hilb}}
\def\Curvesstack{\mathcal{C}\!\mathit{urves}}
\def\Polarizedstack{\mathcal{P}\!\mathit{olarized}}
\def\Complexesstack{\mathcal{C}\!\mathit{omplexes}}
% \Pic is the operator that assigns to X its picard group, usage \Pic(X)
% \Picardstack_{X/B} denotes the Picard stack of X over B
% \Picardfunctor_{X/B} denotes the Picard functor of X over B
\def\Pic{\mathop{\mathrm{Pic}}\nolimits}
\def\Picardstack{\mathcal{P}\!\mathit{ic}}
\def\Picardfunctor{\mathrm{Pic}}
\def\Deformationcategory{\mathcal{D}\!\mathit{ef}}
\usepackage{cleveref}
\crefname{originalremark}{original remark}{original remarks}
\Crefname{originalremark}{Original remark}{Original remarks}
\makeatletter
\newlabel{descent-remark-list-local-source-target}{{04R3}{source}{}{sourcepointer-1}{}}
\newlabel{descent-remark-list-local-source-target@cref}{{[originalremark][0][]04R3}{[1][0][]source}{}{}{}}
\makeatother
\crefname{originallemma}{original lemma}{original lemmas}
\Crefname{originallemma}{Original lemma}{Original lemmas}
\makeatletter
\newlabel{lemma-base-change-quasi-compact}{{03HF}{source}{}{sourcepointer-2}{}}
\newlabel{lemma-base-change-quasi-compact@cref}{{[originallemma][0][]03HF}{[1][0][]source}{}{}{}}
\makeatother
\crefname{originallemma}{original lemma}{original lemmas}
\Crefname{originallemma}{Original lemma}{Original lemmas}
\makeatletter
\newlabel{lemma-local-source-target-at-point}{{04NC}{source}{}{sourcepointer-3}{}}
\newlabel{lemma-local-source-target-at-point@cref}{{[originallemma][0][]04NC}{[1][0][]source}{}{}{}}
\makeatother
\crefname{originallemma}{original lemma}{original lemmas}
\Crefname{originallemma}{Original lemma}{Original lemmas}
\makeatletter
\newlabel{lemma-local-source-target-global-implies-local}{{04RF}{source}{}{sourcepointer-4}{}}
\newlabel{lemma-local-source-target-global-implies-local@cref}{{[originallemma][0][]04RF}{[1][0][]source}{}{}{}}
\makeatother
\crefname{originallemma}{original lemma}{original lemmas}
\Crefname{originallemma}{Original lemma}{Original lemmas}
\makeatletter
\newlabel{morphisms-lemma-base-change-quasi-finite}{{01TM}{source}{}{sourcepointer-5}{}}
\newlabel{morphisms-lemma-base-change-quasi-finite@cref}{{[originallemma][0][]01TM}{[1][0][]source}{}{}{}}
\makeatother
\crefname{originallemma}{original lemma}{original lemmas}
\Crefname{originallemma}{Original lemma}{Original lemmas}
\makeatletter
\newlabel{morphisms-lemma-quasi-finite-points-open}{{01TI}{source}{}{sourcepointer-6}{}}
\newlabel{morphisms-lemma-quasi-finite-points-open@cref}{{[originallemma][0][]01TI}{[1][0][]source}{}{}{}}
\makeatother
\title{The quasi-finite locus is open and commutes with base change}
\author{The Stacks Project Authors}
\date{stacks-project:spaces-morphisms.tex:local-snapshot-20261012; GNU FDL 1.2 or later}
\begin{document}
\maketitle
\setcounter{subsection}{1}
\noindent\textit{Editorial note.} All delivered mathematical text is the author's own, copied from the pinned Stacks Project chapter file spaces-morphisms.tex; no statement or proof step was rewritten, reordered or removed.\par\medskip
\noindent\textit{Editorial note.} The proof is one paragraph reducing all four assertions to the scheme case through a single double presentation; nothing is left to the reader beyond applying the cited scheme statements.\par\medskip
\noindent\textit{Editorial note.} The six references of the counted spans are recorded as non-load-bearing navigation pointers; the argument that consumes them stays inside the counted proof span.\par\medskip
\noindent\textit{Editorial note.} Preamble: the three native excerpt spans of sources/stacks-spaces-morphisms/preamble.tex (lines 8-13, 14-18 and 154-217). The chapter file itself carries no preamble text.\par\medskip
\noindent\textit{Allowed standard prerequisites (editorial).}\begin{itemize}
\item Local quasi-finiteness is etale local on the source-and-target. Descent, Remark descent-remark-list-local-source-target, used at proof lines 5363-5365. Lets the statements be checked on the scheme presentations U -> V and U' -> V'.
\item Pointwise and global forms of the source-and-target locality for this chapter. Lemmas lemma-local-source-target-at-point and lemma-local-source-target-global-implies-local, used at proof lines 5366-5370. Convert the local-on-presentations statements back to f and f'.
\item Openness and base change behaviour of the quasi-finite locus for schemes. Morphisms, Lemmas morphisms-lemma-base-change-quasi-finite and morphisms-lemma-quasi-finite-points-open, used at proof lines 5371-5373. Are the scheme inputs for parts (1), (2) and (3).
\item Quasi-finiteness is local quasi-finiteness plus quasi-compactness, which is stable under base change. Lemma lemma-base-change-quasi-compact, used at proof line 5374. Deduces part (4) from part (3).
\end{itemize}\medskip
\ifdefined\hypertarget\hypertarget{sourcepointer-1}{}\fi
\noindent\textit{Source reference.} Original remark 04R3: focused-seeds/AGP/sources/stacks-snapshot-20261011/tags.txt, source lines 6171--6171. Navigation pointer into the Stacks Project; the standard fact it supplies is applied inside the counted proof span and the argument resting on it stays there.\par\medskip
\ifdefined\hypertarget\hypertarget{sourcepointer-2}{}\fi
\noindent\textit{Source reference.} Original lemma 03HF: focused-seeds/AGP/sources/stacks-snapshot-20261011/tags.txt, source lines 4565--4565. Navigation pointer into the Stacks Project; the standard fact it supplies is applied inside the counted proof span and the argument resting on it stays there.\par\medskip
\ifdefined\hypertarget\hypertarget{sourcepointer-3}{}\fi
\noindent\textit{Source reference.} Original lemma 04NC: focused-seeds/AGP/sources/stacks-snapshot-20261011/tags.txt, source lines 6073--6073. Navigation pointer into the Stacks Project; the standard fact it supplies is applied inside the counted proof span and the argument resting on it stays there.\par\medskip
\ifdefined\hypertarget\hypertarget{sourcepointer-4}{}\fi
\noindent\textit{Source reference.} Original lemma 04RF: focused-seeds/AGP/sources/stacks-snapshot-20261011/tags.txt, source lines 6183--6183. Navigation pointer into the Stacks Project; the standard fact it supplies is applied inside the counted proof span and the argument resting on it stays there.\par\medskip
\ifdefined\hypertarget\hypertarget{sourcepointer-5}{}\fi
\noindent\textit{Source reference.} Original lemma 01TM: focused-seeds/AGP/sources/stacks-snapshot-20261011/tags.txt, source lines 2379--2379. Navigation pointer into the Stacks Project; the standard fact it supplies is applied inside the counted proof span and the argument resting on it stays there.\par\medskip
\ifdefined\hypertarget\hypertarget{sourcepointer-6}{}\fi
\noindent\textit{Source reference.} Original lemma 01TI: focused-seeds/AGP/sources/stacks-snapshot-20261011/tags.txt, source lines 2375--2375. Navigation pointer into the Stacks Project; the standard fact it supplies is applied inside the counted proof span and the argument resting on it stays there.\par\medskip
\begin{lemma}
\label{lemma-base-change-quasi-finite-locus}
Let $S$ be a scheme. Let $f : X \to Y$ and $g : Y' \to Y$ be morphisms
of algebraic spaces over $S$. Denote $f' : X' \to Y'$ the base change
of $f$ by $g$. Denote $g' : X' \to X$ the projection.
Assume $f$ is locally of finite type.
Let $W \subset |X|$, resp.\ $W' \subset |X'|$
be the set of points where $f$, resp.\ $f'$ is quasi-finite.
\begin{enumerate}
\item $W \subset |X|$ and $W' \subset |X'|$ are open,
\item $W' = (g')^{-1}(W)$, i.e., formation of the locus where
$f$ is quasi-finite commutes with base change,
\item the base change of a locally quasi-finite morphism is
locally quasi-finite, and
\item the base change of a quasi-finite morphism is quasi-finite.
\end{enumerate}
\end{lemma}

\begin{proof}
Choose a scheme $V$ and a surjective \'etale morphism $V \to Y$.
Choose a scheme $U$ and a surjective \'etale morphism $U \to V \times_Y X$.
Choose a scheme $V'$ and a surjective \'etale morphism $V' \to Y' \times_Y V$.
Set $U' = V' \times_V U$ so that $U' \to X'$ is a surjective \'etale
morphism as well. Picture
$$
\vcenter{
\xymatrix{
U' \ar[d] \ar[r] & U \ar[d] \\
V' \ar[r] & V
}
}
\quad\text{lying over}\quad
\vcenter{
\xymatrix{
X' \ar[d] \ar[r] & X \ar[d] \\
Y' \ar[r] & Y
}
}
$$
Choose $u \in |U|$ with image $x \in |X|$.
The property of being "locally quasi-finite" is \'etale local on
the source-and-target, see
Descent, Remark \ref{descent-remark-list-local-source-target}.
Hence Lemmas \ref{lemma-local-source-target-at-point} and
\ref{lemma-local-source-target-global-implies-local} apply and
we see that $f : X \to Y$ is quasi-finite at $x$
if and only if $U \to V$ is quasi-finite at $u$.
Similarly for $f' : X' \to Y'$ and the morphism $U' \to V'$.
Hence parts (1), (2), and (3) reduce to
Morphisms, Lemmas \ref{morphisms-lemma-base-change-quasi-finite} and
\ref{morphisms-lemma-quasi-finite-points-open}.
Part (4) follows from (3) and Lemma \ref{lemma-base-change-quasi-compact}.
\end{proof}

\end{document}
